\section{保持原始支持的精确编译}\label{sec:compile}
局部化提供一个足够小的候选接口，但它尚未回答哪些候选组合合法。本节给出从候选表示到完整支持的构造合同。一般正则关系监视、图消元、选择分组和支持编码已有成熟基础~\cite{kater2023,kucera2021}；本文对这些步骤的使用，必须保证不会丢失前端特意保留的来源身份与 CO 相关性。

\subsection{稀疏生成 CO 与联合关系监视}
令 $g$ 包含 $\Cov(L)$ 和每个未决同址写对所选的一条方向。若 $g$ 无环，则 $g^+$ 恰为扩展 $L$ 的严格全序：有限偏序的覆盖闭包等于自身，而所有未决写对已经选择了一个方向，所以每对写可比；无环性保证严格性。反之，对任意扩展 $L$ 的 CO 总序按其方向选择未决对，得到无环 $g$ 且 $g^+=\CO$。这证明稀疏生成没有预先排除合法写序。已固定的传递比较依然作为原始常量标签存在。

把相邻程序序和初始化生成边记为字母 $s$，RF 记为 $r$，逆 RF 记为 $\bar r$，CO 生成边 $g$ 记为 $c$。式~\eqref{eq:ra} 的禁止闭合游走可用正则语言
\begin{equation}
 \mathcal L=(s\mid r)^+(r\mid c\mid\bar r c^+)^*\label{eq:lang}
\end{equation}
检查。第一段为非空 HB，第二段允许空的 ECO；$c^+$ 实现 CO 闭包，$\bar r c^+$ 实现 FR。该语言有四状态部分 DFA；另用单状态无环监视器检查 $g$。读来源恰一约束由局部计数或原子来源选择实施。每次选择加入\emph{真实}关系边，而不把识别器的所有 $H$ 事实当作已选 RF。

为避免遗忘禁止游走的起点，在 RA 乘积图中为每个事件增加从接受状态回到初态的\emph{标记边}。令普通乘积路径关系为 $R=G^*$，标记边集为 $B_m$，并保存
\[
 T=R;B_m;R.
\]
$T$ 表示恰经过一个标记的路径。原图存在 $\mathcal L$ 闭合游走，当且仅当 $T$ 对角非空：接受路径接一条标记边即闭合，反向将唯一标记旋转到末尾即可。不得对 $T$ 再作传递闭包，因为经过两个标记的环不一定表示一个被禁止的词。

在组 $j$ 处理完后，静态后缀候选集合包含所有剩余替代项可能增加的普通边及尚未加入的标记边。令 $I_j$ 为已出现且仍可能接收后缀边的乘积顶点，$O_j$ 为仍可能发出后缀边的乘积顶点。摘要只保留
\begin{equation}
 \left(R|_{I_j\times O_j},\ T|_{I_j\times O_j}\right),\label{eq:ports}
\end{equation}
并与 CO 监视器、局部度计数组成一个\emph{联合}状态。后缀候选可能互不相容，它们仅决定保留坐标，不能作为实际边加入任何历史。

\begin{clemma}[联合端口的后缀等价]\label{lem:ports}
两个未被拒绝的同层前缀若具有相同联合摘要和计数，则任意相同的后缀选择在二者上同时合法或同时非法。
\end{clemma}
\begin{cproof}
新产生的禁止闭合游走必含后缀新边，否则此前已经拒绝。将游走旋转到一条新边处，并切成新边与历史最大连续片段。每个历史片段从新边目标进入、从新边来源离开，端点必在 $I_j\times O_j$。恰一标记的游走中，至多一个历史片段使用 $T$，其余使用 $R$；若标记来自新边，则所有历史片段均使用 $R$。相同摘要允许在另一历史内替换每段同端点、同标记数路径，得到相同违反；反向同理。这里允许路径重复顶点，符合闭合游走语义。局部恰一条件因计数与剩余组相同而同判。
\end{cproof}

历史片段的替换始终在\emph{同一个}历史内完成。分别合并不同监视器的可达关系，或将互不相容历史的边取并，会失去这个论证。附录~\ref{app:update} 给出闭包更新公式及两个标记不能组合的反例。

\subsection{标签、支持扫描与 UP}
编译器逐组展开选择，用引理~\ref{lem:ports} 合并同层等价摘要，得到有向无环图 $\mathcal G$。每条边携带该组对原始变量的赋值标签；每个原始位沿完整路径恰被赋值一次。被删 RF、单例 RF 和已固定 CO 在根部常量标签中赋值。不同前缀即使合并为同一状态，其入路径标签也保持不同。

逐层归纳可知，接受路径的原始投影恰为 $\mathcal V$。给定 $\alpha$，禁用含相反文字的标签边，分别进行根到节点的前向扫描和节点到接受端的反向扫描。一条未禁用边有支持，当且仅当其来源可从根到达且其目标可到达接受端；汇总这些边的标签得到 $\supp(\alpha)$。不存在接受路径时报告 UNSAT。每次查询成本线性于显式图、标签以及输入／输出长度；撤销假设只需重新扫描。

同一过程可以编码为支持 CNF~\cite{kucera2021}。实现使用其对偶形式：辅助位表示标签被禁用、节点无法前向到达／无法后向接受，以及边已经失去路径支持。原始假设触发相反标签失效；按拓扑序，若所有入边都失效则推出该节点无法前向到达，反向同理。边的标签被禁用、来源不可达或目标不能接受，都足以推出该边失效。对每个原始文字 $\lambda$，加入“含 $\lambda$ 的所有边都失效则 $\neg\lambda$”。UP 因而推出所有无支持文字的否定；若根不能到达接受端，则与根必须可接受的约束冲突。存在合法路径时，将这些辅助位取为其真实失效状态给出满足扩展，故不会产生不可靠推导。这是原始变量的域一致性合同，不是关于任意辅助赋值的更强结论。

\subsection{构造成本与原子 RF 分组}
宽度 $k$ 时，RA 乘积最多有 $4k$ 个活跃顶点，其 $R,T$ 两矩阵合计至多 $32k^2$ 位；CO 普通可达矩阵再用 $k^2$ 位，来源度计数至多 $k$ 个布尔位。联合状态数因而至多 $S=2^{33k^2+k}$。简单候选图有 $O(kN)$ 条边；二元选择组及固定边组数量同阶。每组最多两项，生成图、标签和支持 CNF 得式~\eqref{eq:size}；矩阵更新带来 $\operatorname{poly}(k)$ 因子。

识别阶段中 $D$ 只减，$M,B_\forall,H,L$ 只增，总计至多 $O(N^2)$ 次严格变化。显式矩阵实现每轮保守计 $O(I^3)$，于是粗界 $O(I^5)$ 足够覆盖识别。该界不是精细运行时间预测；保存完整逐条证明轨迹时，轨迹空间需另外计费。至此完成定理~\ref{thm:main}(D) 的证明。

实际实现将一个读的 $t_r=|D(r)|$ 个 RF 二元位合为一个 $t_r$ 路原子组。选择第 $j$ 个来源时，标签同时将该 RF 位置真、其余 $t_r-1$ 位置假。它保留来源身份，只消除了不满足恰一条件的局部组合。

\begin{cprop}[原子分组的保守界]\label{prop:atomic}
定理~\ref{thm:main} 的输入采用原子 RF 组时，令
\[
 k_A=(2\delta+1)b+a,\qquad S_A=2^{33k_A^2+k_A}.
\]
表示大小为 $O(|X|+k_A^2NS_A)$，原始支持合同不变。
\end{cprop}
\begin{cproof}
来源各距读至多 $\delta$，但两个来源之间可能相距 $2\delta$。因此超组跨度应按 $2\delta$ 计；按最大端点排序后的存活窗口给出 $k_A$，不能直接沿用二元界 $k$。每个 $t_r$ 路组有 $t_r$ 个长度 $t_r$ 的标签，标签成本是 $t_r^2$。由 $t_r\le k_A$ 得 $\sum_r t_r^2\le Nk_A^2$；CO 二路组仅需 $O(k_AN)$，按每层状态界复制即得结论。每个合法 RF 恰一赋值与一个组分支一一对应，完整生成边集合相同；非法前缀不会因继续加边而变合法，所以重排选择组不丢合法模型。
\end{cproof}

后处理采用标准的同层、同标签后缀语言商并删除无接受后缀的节点。它减小查询图，但不能降低已经发生的构造峰值。可选 C++ 位集实现执行同一更新代数，产品接口超过 128 位时回退 Python；它不改变摘要等价关系。本文分别报告状态减少、后处理压缩和底层执行成本，避免把它们都归于局部化定理。
